\chapter{Méthodes de pénalité et Lagrangien augmenté}
\label{chap:penalty-augmented-lagrangian}

\section{La méthode de pénalité quadratique}

\subsection{Motivation}

\begin{algorithm}[H]
    \caption{Quadratic penalty method}
    \SetAlgoLined
    Given $\mu_0>0$, a nonnegative sequence $(\tau_k)_k$ with $\tau_k \to 0$, and a starting point $x_0^s$\;
    \For{$k=0, 1, \ldots$}{
        Find an approximate minimizer $x_k$ of $Q(\cdot;\mu_k)$, starting at $x_k^s$ and terminating when $\|\nabla_x Q(x;\mu_k)\|\leq \tau_k$\;
        \If{final convergence test satisfied}{
            \Stop with approximate solution $x_k$\;
        }
        Choose new penalty parameter $\mu_{k+1}>\mu_k$\;
        Choose new starting point $x_{k+1}^s$\;
    }
\end{algorithm}

\subsection{Convergence de la méthode (avec contraintes d'égalité)}

% TODO: Contenu à développer

\section{Fonctions de pénalité non lisses}

\subsection{Propriétés}

\begin{algorithm}[H]
    \caption{Classical $\ell^1$ penalty method}
    \SetAlgoLined
    Given $\mu_0>0$, tolerance $\tau>0$, starting point $x_0^s$\;
    \For{$k=0, 1, \ldots$}{
        Find an approximate minimizer $x_k$ of $\phi_1(x;\mu_k)$, starting at $x_k^s$\;
        \If{$\|h(x_k)\|\leq \tau$}{
            \Stop with approximate solution $x_k$\;
        }
        Choose new penalty parameter $\mu_{k+1}>\mu_k$\;
        Choose new starting point $x_{k+1}^s$\;
    }
\end{algorithm}

\subsection{Méthode de pénalité $\ell^1$ pratique}

% TODO: Contenu à développer

\section{Méthode du Lagrangien augmenté}

\subsection{Propriétés du cas avec contraintes d'égalité}

\begin{algorithm}[H]
    \caption{Augmented Lagrangian method for equality constraints}
    \SetAlgoLined
    Given $\mu_0>0$, tolerance $\tau_0>0$, starting points $x_0^s$ and $\lambda^0$\;
    \For{$k=0, 1, \ldots$}{
        Find an approximate minimizer $x_k$ of $\mathcal{L}_A(\cdot, \lambda^k;\mu_k)$, starting at $x_k^s$, and terminating when $\|\nabla_x \mathcal{L}_A(x_k, \lambda^k; \mu_k)\| \leq \tau_k$\;
        \If{a convergence test for the equality-constrained problem is satisfied}{
            \Stop with approximate solution $x_k$\;
        }
        Update Lagrange multipliers using
        \[
            \lambda_i^{k+1} = \lambda_i^k - \mu_k c_i(x_k), \text{ for all } i\in\mathcal{E};
        \]
        Choose new penalty parameter $\mu_{k+1} \geq \mu_k$\;
        Set starting point for the next iteration to $x_{k+1}^s = x_k$\;
        Select tolerance $\tau_{k+1}$\;
    }
\end{algorithm}

\subsection{Méthodes pratiques du Lagrangien augmenté}

\subsubsection{Formulation avec contraintes de bornes}

\begin{algorithm}[H]
    \caption{Bound-constrained Lagrangian method}
    \SetAlgoLined
    Choose an initial point $x_0$ and initial multipliers $\lambda^0$\;
    Choose convergence tolerances $\eta_*$ and $\omega_*$\;
    Set $\mu_0=10$, $\omega_0=\frac{1}{\mu_0}$, and $\eta_0=\frac{1}{\mu_0^{0.1}}$\;
    \For{$k=0,1, \ldots$}{
        Find an approximate solution $x_k$ of the problem
        \[
            \min_x \mathcal{L}_A(x,\lambda^k;\mu_k) \text{ subject to } l\leq x\leq u
        \]
        such that
        \[
            \left\| x_k-P\left(x_k - \nabla_x \mathcal{L}_A(x_k, \lambda^k ; \mu_k), l, u\right) \right\| \leq \omega_k;
        \]
        \If{$\|c(x_k)\|\leq \eta_k$}{
            \Comment{test for convergence}
            \If{$\|c(x_k)\|\leq \eta_*$ \And $\left\| x_k-P\left(x_k - \nabla_x \mathcal{L}_A(x_k, \lambda^k ; \mu_k), l, u\right) \right\| \leq \omega_*$}{
                \Stop with approximate solution $x_k$\;
            }
            \Comment{update multipliers, tighten tolerances}
            $\lambda^{k+1}\leftarrow \lambda^k - \mu_k c(x_k)$\;
            $\mu_{k+1}\leftarrow \mu_k$\;
            $\eta_{k+1}\leftarrow \frac{\eta_k}{\mu_{k+1}^{0.9}}$\;
            $\omega_{k+1} \leftarrow \frac{\omega_k}{\mu_{k+1}}$\;
        }
        \Else{
            \Comment{increase penalty parameters, tighten tolerances}
            $\lambda^{k+1}\leftarrow \lambda^k$\;
            $\mu_{k+1}\leftarrow 100 \mu_k$\;
            $\eta_{k+1}\leftarrow \frac{1}{\mu_{k+1}^{0.1}}$\;
            $\omega_{k+1} \leftarrow \frac{1}{\mu_{k+1}}$\;
        }
    }
\end{algorithm}

\subsubsection{Formulation avec contraintes linéaires}

% TODO: Contenu à développer

\subsubsection{Formulation sans contraintes}

% TODO: Contenu à développer
